\documentclass{ExcelAtFIT}

\ExcelFinalCopy

\hypersetup{
	pdftitle={3D Neural Cellular Automata: An Interactive Framework for Volumetric Morphogenesis},
	pdfauthor={Michal Repčík},
	pdfkeywords={Neural Cellular Automata, Blender, Voxel, Morphogenesis, Deep Learning}
}

% Poster figure mapping (figures are on the poster; not embedded here):
% Fig. 1: Rendered mesh + voxelized target in Blender (mesh -> dataset creation and export)
% Fig. 2: Blender viewport (mesh, voxel target, livestreamed snapshots) + UI panels/settings
% Fig. 3: One training step of a 3D NCA
% Fig. 4: Metamorphogenesis loss curve (framework visualization utilities)
% Fig. 5: Metamorphogenesis training targets (framework visualization utilities)

\ExcelYear{2026}

\PaperTitle{3D Neural Cellular Automata: An Interactive Framework for Volumetric Morphogenesis}

\Authors{Michal Repcik*}
\affiliation{*%
  \href{mailto:xrepcim00@stud.fit.vutbr.cz}{xrepcim00@stud.fit.vutbr.cz},
  \textit{Faculty of Information Technology, Brno University of Technology}}

\Abstract{
Existing 3D Neural Cellular Automata implementations are mostly batch scripts with no
interactive training tooling and no practical workflow for modelling targets and exporting ready-to-use
coloured voxel datasets (Fig.~1).
%
This work presents a modular framework consisting of a Blender add-on for mesh
voxelization and live viewport rendering (Fig.~2), a TCP server exposing any training backend
over a local or tunnelled connection, and a standalone Python core package usable in Jupyter
notebooks without Blender; one training step is summarised in Fig.~3.
%
The framework is validated through morphogenesis experiments on symmetric and
asymmetric targets, and through multi-target metamorphogenesis, where a single model
grows distinct voxel shapes conditioned on task channels (Figs.~4--5).
%
Implementing a new NCA experiment requires subclassing one abstract class; the
server, protocol, logging, and visualisation pipeline require no changes.\footnote{Demo video: \url{https://youtu.be/hMkapi_rayU}.}
}

\begin{document}

\startdocument

%%%%%%%%%%%%%%%%%%%%%%%%%%%%%%%%%%%%%%%%%%%%%%%%%%%%%%%%%%%%
\section{Introduction}

Neural Cellular Automata (NCA)~\cite{Mordvintsev2020} replace the hand-crafted rules
of classical cellular automata with a small shared neural network trained by gradient
descent. Each voxel perceives its local neighbourhood through fixed convolutional
filters and applies a learned residual update. Complex global structure -- growth,
regeneration, pattern formation -- emerges from repeated local interaction without
any centralised coordination.

Extending NCA to three-dimensional voxel grids is conceptually straightforward.
Prior work confirms that volumetric growth is feasible~\cite{Sudhakaran2021}\footnote{Framework source code (E. Rossi 2021 base): \url{https://github.com/e-rossi/3d-nca}.}\footnote{Extension repo (A. Jha 2024): \url{https://github.com/Aadityaza/3d-Growing-neural-cellular-automata}.}, but most implementations remain standalone training scripts.
They lack live visualisation, interactive hyperparameter steering, and a mesh-to-voxel dataset workflow.

This work presents a framework that addresses these limitations. It consists of three
decoupled layers: a Blender add-on, a TCP training server, and a standalone core
package. The central design goal is that implementing any NCA experiment requires
writing one class; everything else is provided by the framework.

Main contributions are: (1) a mesh-to-voxel dataset workflow integrated into Blender, including colour-preserving
export of ready-to-use voxel targets (Fig.~1); (2) live 3D training visualisation and interactive control through
the Blender UI (Fig.~2) with a decoupled backend boundary; and (3) model-level extensions --- positional channels and
task conditioning --- that address asymmetric targets and enable multi-target metamorphogenesis (Figs.~4--5).

%%%%%%%%%%%%%%%%%%%%%%%%%%%%%%%%%%%%%%%%%%%%%%%%%%%%%%%%%%%%
\section{System Architecture}

\textbf{Blender add-on.}
The add-on performs mesh voxelization and dataset export (Fig.~1) and viewport rendering of live training snapshots
(Fig.~2). A selected mesh is sampled onto a
3D grid using a raycast inside-outside test. Per-voxel colour is extracted either from the Principled BSDF
base colour (fast) or via UV texture lookup (general). The voxelized target can be exported as \texttt{.npz}
or transmitted directly to the training server.

During training, state broadcasts received from the server are rendered as
merged-cube meshes in the Blender viewport. The source mesh, voxelized target, and
current NCA state appear side by side in three spatial slots for direct visual
comparison during training.

The sidebar control panel exposes connection settings, session control, and a schedule editor for
one-shot parameter updates (Fig.~2).

\textbf{Server and protocol.}
The server exposes training over TCP and runs the selected backend in a background thread while
streaming state snapshots at a controlled cadence. For remote GPU training on Colab or Kaggle, an ngrok
TCP tunnel bridges the connection with no firewall configuration required.

\textbf{Core package.}
The core package provides a clean PyTorch NCA implementation with configuration dataclasses and
self-describing checkpoints. It is independently usable from Python scripts or Jupyter notebooks via the
high-level \texttt{NCAModel} API.

%%%%%%%%%%%%%%%%%%%%%%%%%%%%%%%%%%%%%%%%%%%%%%%%%%%%%%%%%%%%
\section{NCA Model}

\textbf{State representation.}
Each voxel holds a real-valued vector of $C$ channels partitioned into hidden,
positional, task, and visible groups. Hidden channels carry unsupervised internal
memory. Visible channels (RGBA) define the output compared against the target.
Positional channels encode normalised absolute coordinates
$(x,y,z) \in [0,1]^3$, re-injected as read-only constants after each step to
provide global spatial reference that purely local symmetric perception cannot
supply. Task channels enable one model to learn multiple target behaviours via
one-hot conditioning.

\textbf{Perception and update.}
Fixed $3\times3\times3$ depthwise convolutions produce a perception vector at every
cell. The baseline configuration uses three filter groups: identity, 6-neighbour
sum, and Laplacian. The directional configuration uses five groups: identity,
6-neighbour sum, and signed gradients along $x$, $y$, and $z$. Directional filters
resolve reflection ambiguity for asymmetric targets, where the isotropic baseline
converges to mirrored solutions.

Combined with positional channels $(x,y,z)$ as read-only constants, the update rule receives an explicit global
frame of reference while remaining strictly local and translation-invariant in its learned parameters.

The update MLP maps the perception vector to a bounded state increment,
$\Delta \mathbf{s} = 0.1 \cdot \tanh(\mathrm{MLP}(\mathrm{perceived}))$,
applied residually: $\mathbf{s}_{t+1} = \mathbf{s}_t + \Delta \mathbf{s}$.
An alive mask computed by $3\times3\times3$ max-pooling of the alpha channel and
thresholding gates which cells receive updates. Dead cells adjacent to alive cells
are included in the mask and can grow; isolated dead regions cannot activate.

\textbf{Training.}
Pool-based training maintains a pool of persistent states across developmental stages.
Each iteration runs the automaton for a randomly sampled number of steps under a step-range curriculum,
reseeds the highest-loss sample to maintain diversity, and optimises a weighted sum of alpha loss,
masked colour loss, and an overflow penalty for activity outside the target volume.
In the default configuration, the pool holds 32 states and the step range expands from 8--16 steps early in training
to 32--64 steps later to encourage stability over longer horizons.

%%%%%%%%%%%%%%%%%%%%%%%%%%%%%%%%%%%%%%%%%%%%%%%%%%%%%%%%%%%%
\section{Experiments}

All experiments use a $32^3$ grid, 16 hidden channels, and 4 visible (RGBA) channels,
seeded from a single alive centre cell.

\textbf{Morphogenesis.}
A symmetric target converged from total loss 2.36 to $4.5 \times 10^{-3}$ over
3000 epochs. An asymmetric target (chess pawn) with the symmetric baseline
configuration repeatedly converged to mirrored intermediate solutions. Enabling
positional channels and directional perception resolved this; the corrected run
reached a total loss of $2.75 \times 10^{-4}$ at epoch 4000. This failure and its
resolution directly motivated positional channels as a framework-level feature.

\textbf{Metamorphogenesis.}
Metamorphogenesis denotes a multi-target regime in which one model both generates form and transitions
between forms through task-channel conditioning (Figs.~4--5). A related approach to conditional morphogenesis
in NCA was published
concurrently~\cite{Sakour2025}; the implementation here demonstrates that the same
conditioning idea is achievable with minimal code changes to the framework, a custom
\texttt{NCARunner} subclass with no modifications to the server, protocol, or
visualisation pipeline, confirming the practical extensibility of the design.

Training uses a phased curriculum: the model first learns a single target, then expands to additional
targets at phase boundaries (Fig.~5). Pool states can be relabelled across phases to preserve continuity of the
transition, producing characteristic loss spikes followed by re-convergence (Fig.~4).

%%%%%%%%%%%%%%%%%%%%%%%%%%%%%%%%%%%%%%%%%%%%%%%%%%%%%%%%%%%%
\section{Conclusions}

This work presents a modular 3D NCA framework connecting mesh-based dataset preparation, remote GPU
training, and live Blender viewport visualisation in a single interactive workflow.
Experiments validate symmetric/asymmetric morphogenesis and task-conditioned metamorphogenesis,
and show positional channels as an effective remedy for asymmetric reflection ambiguity.

\section*{Acknowledgements}
I would like to thank my supervisor for guidance throughout this work.

\phantomsection
\bibliographystyle{unsrt}
\bibliography{bibliography}

\end{document}
